% Generated draft: requires mathematical and source review.
\documentclass[11pt]{article}
\usepackage[margin=25mm]{geometry}
\usepackage{fontspec}
\setmainfont{Latin Modern Roman}
\usepackage{amsmath,amssymb,mathtools,unicode-math}
\setmathfont{Latin Modern Math}
\usepackage{enumitem,needspace,xurl,hyperref}
\setlength{\parindent}{0pt}
\setlength{\parskip}{5pt}
\setlength{\emergencystretch}{3em}
\newcommand{\N}{\mathbb{N}}
\newcommand{\Z}{\mathbb{Z}}
\newcommand{\Q}{\mathbb{Q}}
\newcommand{\R}{\mathbb{R}}
\newcommand{\C}{\mathbb{C}}
\newcommand{\F}{\mathbb{F}}
\newcommand{\A}{\mathbb{A}}
\newcommand{\PP}{\mathbb{P}}
\newcommand{\E}{\mathbb{E}}
\newcommand{\eps}{\varepsilon}
\newcommand{\dd}{\,\mathrm{d}}
\begin{document}
\section*{UnsolvedMath ID: 800804}
\textbf{Source problem code:} catalogue entry
\subsection{Definitions and mathematical statement}
The supplied record asks for a geometric description, rather than asserting a particular curvature formula or existence theorem. Write \(M\) for the underlying manifold and \(\mu\) for its volume structure. The comparison object in the source is the group of volume-preserving diffeomorphisms,
\[
\mathscr{D}_{\mu}(M)=\{\phi\in\operatorname{Diff}(M):\phi^*\mu=\mu\},
\]
with the regularity and boundary conventions left unspecified by the supplied record.

Let \(\mathcal G\rightrightarrows B\) denote the intended Lie groupoid of multiphase or generalized volume-preserving diffeomorphisms, with source and target maps \(s,t\). This notation is conditional: the supplied record does not define \(\mathcal G\), its objects, or its arrows. A source-wise metric is a family of metrics \(G_b\) on the source fibers \(s^{-1}(b)\), for \(b\in B\). With the convention that \(kh\) is defined when \(s(k)=t(h)\), right multiplication by an arrow \(h\) is
\[
R_h:s^{-1}(t(h))\longrightarrow s^{-1}(s(h)),\qquad k\longmapsto kh.
\]
Right invariance means
\[
R_h^*G_{s(h)}=G_{t(h)}\qquad\text{for every }h\in\mathcal G.
\]
The requested metric is the source's \(L^2\) metric. Its integral formula cannot be specified from this record: the phase data, velocity representation, integration measures, and underlying pointwise inner products are not supplied.

\textbf{Problem.} Describe the differential geometry of this right-invariant, source-wise \(L^2\) metric on the Lie groupoid of multiphase or generalized volume-preserving diffeomorphisms, analogously to Arnold's description for \(\mathscr{D}_{\mu}(M)\). In particular, compute sectional curvatures and, wherever a mathematically defined Ricci contraction exists, Ricci curvatures, and describe conjugate points along geodesics in the source fibers.

For precision concerning these requested quantities, whenever a source fiber admits the requisite Levi--Civita connection \(\nabla\), use the convention
\[
R(X,Y)Z=\nabla_X\nabla_YZ-\nabla_Y\nabla_XZ-\nabla_{[X,Y]}Z.
\]
For linearly independent tangent vectors \(u,v\) at a point of that fiber, its sectional curvature is
\[
K(u,v)=\frac{\langle R(u,v)v,u\rangle}
{\langle u,u\rangle\langle v,v\rangle-\langle u,v\rangle^2}.
\]
The formal Ricci contraction is
\[
\operatorname{Ric}(u,v)=\operatorname{tr}\bigl(w\longmapsto R(w,u)v\bigr),
\]
but an infinite-dimensional setting requires a justified notion and existence of this trace; neither is supplied. Along a geodesic \(\gamma\), a time \(T>0\) is conjugate to time \(0\) in the Jacobi-field sense if there is a nonzero Jacobi field \(J\) satisfying
\[
\nabla_{\dot\gamma}^{\,2}J+R(J,\dot\gamma)\dot\gamma=0,
\qquad J(0)=J(T)=0.
\]
These definitions specify the requested geometric quantities where they exist; they do not assert that the necessary analytic structures exist for the intended groupoid. A fully specified formulation still requires the source's definitions of multiphase and generalized diffeomorphisms, its \(L^2\) metric, and its manifold, regularity, and boundary hypotheses.
\subsection{Short English statement}
Describe the geometry of the source-wise, right-invariant \(L^2\) metric on the intended groupoid of multiphase or generalized volume-preserving diffeomorphisms, including sectional and Ricci curvature and conjugate points. The supplied record omits the groupoid and metric constructions and the analytic hypotheses needed to make every requested quantity precise.
\subsection{Sources}
\begin{itemize}
\item Local UnsolvedMath registry, permanent ID \(800804\), exact source problem code \texttt{catalogue entry}, title ``Geometric Hydrodynamics in Open Problems --- Problem 3.'' The supplied registry attributes the problem to Khesin, Misiolek, and Modin, with proposed year 2022, and labels it \texttt{open}. This is a supplied status label, not an independently verified mathematical or literature assessment.
\item Khesin, Misiolek, and Modin, \emph{Geometric Hydrodynamics in Open Problems} (2022), Problem 3. Supplied source URL: \url{https://arxiv.org/abs/2205.01143}. The provenance records access on 2026-10-16 and extraction by \texttt{arxiv\_latex}; these are supplied provenance claims, not a claim of access or independent verification in this draft.
\item The statement invokes Arnold's description of the geometry of \(\mathscr{D}_{\mu}(M)\) as its comparison. No specific Arnold reference is supplied.
\end{itemize}
\end{document}
